\begin{PSbox}[unbreakable]{obj}{اهداف یادگیری}

پس از تکمیل این ماژول، دانشجویان قادر خواهند بود:

\begin{enumerate}
\def\labelenumi{\arabic{enumi}.}
\tightlist
\item
  تبدیل‌های خطی را روی متغیرهای تصادفی اعمال کنند
\item
  \lr{\mbox{PDF}} تبدیل‌های غیرخطی را با روش‌های \lr{\mbox{CDF}} و ژاکوبی استخراج کنند
\item
  تبدیل‌ها را در مسایل مهندسی (توان، \lr{\mbox{SNR}}، پوش) به کار ببرند
\item
  نتایج تحلیلی را با شبیه‌سازی \lr{\mbox{MATLAB}} بررسی کنند
\end{enumerate}

\end{PSbox}

\PSsection{5.1}{مقدمه: چرا متغیرهای تصادفی را تبدیل کنیم؟}

در مهندسی، اغلب آماره‌های یک کمیت را می‌دانیم اما به آماره‌های یک کمیت \textbf{مشتق‌شده} نیاز داریم:

{\def\LTcaptype{none} % do not increment counter
\begin{longtable}[c]{@{}ccc@{}}
\toprule\noalign{}
\rowcolor{PSnavy}\PSth{معلوم} & \PSth{موردنیاز} & \PSth{تبدیل} \\
\midrule\noalign{}
\endhead
\bottomrule\noalign{}
\endlastfoot
ولتاژ \lr{\mbox{V}} & توان \lr{\mbox{P}} & \lr{\mbox{P = V\textsuperscript{2}/R}} \\
\lr{\mbox{SNR}} خطی & \lr{\mbox{SNR}} بر حسب \lr{\mbox{dB}} & \lr{\mbox{Y = 10\ensuremath{\cdot}log\textsubscript{10}(X)}} \\
گاوسی \lr{\mbox{I}}، \lr{\mbox{Q}} & پوش & \lr{R = \ensuremath{\surd}(I\textsuperscript{2} + Q\textsuperscript{2})} \\
جریان \lr{\mbox{I}} & توان & \lr{\mbox{P = I\textsuperscript{2}R}} \\
فاصله \lr{\mbox{d}} & تلفات مسیر & \lr{\mbox{L = d\textsuperscript{\ensuremath{\alpha}}}} \\
\end{longtable}
}

\PSsection{5.2}{تبدیل‌های خطی: \lr{\mbox{Y = aX + b}}}

\PSsubsection{نتیجه}

اگر \lr{\mbox{X}} دارای \lr{\mbox{PDF f\textsubscript{X}(x)}} باشد، آنگاه \lr{\mbox{Y = aX + b}} دارای:

\PSdisplay{f_Y(y) = \frac{1}{|a|} f_X\left(\frac{y-b}{a}\right)}

\PSsubsection{میانگین و واریانس}

\begin{itemize}
\tightlist
\item
  \lr{\mbox{E[Y] = aE[X] + b}}
\item
  \lr{\mbox{Var(Y) = a\textsuperscript{2}Var(X)}}
\end{itemize}

\begin{PSbox}[unbreakable]{ex}{مثال مهندسی: تقویت‌کننده}

سیگنال ورودی \lr{X \textasciitilde{} N(0, \ensuremath{\sigma}\textsuperscript{2})} از تقویت‌کننده‌ای با بهره \lr{\mbox{G = 10}} و آفست \lr{\mbox{DC}} یعنی \lr{\mbox{V\textsubscript{0} = 1.5V}} می‌گذرد.

خروجی: \lr{\mbox{Y = 10X + 1.5}}

\begin{itemize}
\tightlist
\item
  \lr{Y \textasciitilde{} N(10\ensuremath{\cdot}0 + 1.5, 10\textsuperscript{2}\ensuremath{\cdot}\ensuremath{\sigma}\textsuperscript{2}) = N(1.5, 100\ensuremath{\sigma}\textsuperscript{2})}
\item
  توان نویز خروجی به اندازه \lr{\mbox{G\textsuperscript{2} = 100}} افزایش یافته است
\end{itemize}

\end{PSbox}

\begin{PSbox}[unbreakable]{ex}{مثال مهندسی: حسگر دما}

خروجی حسگر \lr{\mbox{V = 0.01\ensuremath{\cdot}T + 0.5}} (\lr{\mbox{V}} بر حسب ولت، \lr{\mbox{T}} بر حسب \lr{\mbox{\ensuremath{^{\circ}}C}})

اگر \lr{\mbox{T \textasciitilde{} N(25, 4)}}: \lr{V \textasciitilde{} N(0.01\ensuremath{\times}25 + 0.5, 0.01\textsuperscript{2}\ensuremath{\times}4) = N(0.75, 0.0004)}

\end{PSbox}

\PSsection{5.3}{تبدیل‌های یکنوای عمومی}

\PSsubsection{روش \lr{\mbox{CDF}} (کلی‌ترین روش)}

برای \lr{\mbox{Y = g(X)}}، \lr{\mbox{f\textsubscript{Y}(y)}} را به این ترتیب بیابید:

\begin{enumerate}
\def\labelenumi{\arabic{enumi}.}
\tightlist
\item
  بنویسید \lr{F\textsubscript{Y}(y) = P(Y \ensuremath{\le} y) = P(g(X) \ensuremath{\le} y)}
\item
  نامساوی را برای \lr{\mbox{X}} حل کنید
\item
  بر حسب \lr{\mbox{F\textsubscript{X}}} بیان کنید
\item
  مشتق بگیرید تا \lr{\mbox{f\textsubscript{Y}(y)}} به دست آید
\end{enumerate}

\PSsubsection{فرمول برای \lr{\mbox{g}} یکنوا}

اگر \lr{\mbox{g}} اکیداً صعودی با وارون \lr{\mbox{g\textsuperscript{-1}}} باشد:\PSdisplay{f_Y(y) = f_X(g^{-1}(y)) \cdot \frac{d}{dy}[g^{-1}(y)]}

اگر \lr{\mbox{g}} اکیداً نزولی باشد:\PSdisplay{f_Y(y) = f_X(g^{-1}(y)) \cdot \left|\frac{d}{dy}[g^{-1}(y)]\right|}

ترکیبی (برای هر دو کار می‌کند):\PSdisplay{f_Y(y) = \frac{f_X(x)}{|g'(x)|}\bigg|_{x=g^{-1}(y)}}

\begin{PSbox}[unbreakable]{ex}{مثال مهندسی: \lr{\mbox{SNR}} بر حسب \lr{\mbox{dB}}}

\lr{\mbox{SNR}} خطی: \lr{\mbox{X \textasciitilde{} Exponential(1)}} (توان محوشدگی رایلی نرمال‌شده).\PSbr{}\lr{\mbox{SNR}} بر حسب \lr{\mbox{dB}}: \lr{\mbox{Y = 10\ensuremath{\cdot}log\textsubscript{10}(X)}}

\begin{itemize}
\tightlist
\item
  \lr{\mbox{g(x) = 10\ensuremath{\cdot}log\textsubscript{10}(x)}}، \lr{\mbox{g’(x) = 10/(x\ensuremath{\cdot}ln10)}}
\item
  \lr{\mbox{g\textsuperscript{-1}(y) = 10\textsuperscript{y/10}}}
\end{itemize}

\PSdisplay{f_Y(y) = \frac{f_X(10^{y/10})}{10/(10^{y/10} \cdot \ln 10)} = \frac{\ln 10}{10} \cdot 10^{y/10} \cdot e^{-10^{y/10}}}

این توزیع گاوسی نیست \lr{—} دنباله چپ بلندتری دارد (فروافت‌های عمیق).

\end{PSbox}

\PSsection{5.4}{تبدیل‌های غیرخطی: حالت غیریکنوا}

\PSsubsection{هنگامی که \lr{\mbox{Y = g(X)}} یک‌به‌یک نیست}

اگر چند مقدار \lr{\mbox{x}} به همان \lr{\mbox{y}} نگاشت شوند، روی تمام جواب‌ها جمع کنید:

\PSdisplay{f_Y(y) = \sum_{i} \frac{f_X(x_i)}{|g'(x_i)|}}

که در آن \lr{\mbox{x\textsubscript{1}}}، \lr{\mbox{x\textsubscript{2}}}، \lr{…} تمام ریشه‌های \lr{\mbox{g(x) = y}} هستند.

\begin{PSbox}[unbreakable]{ex}{مثال مهندسی: توان از ولتاژ \lr{\mbox{(Y = X\textsuperscript{2})}}}

ولتاژ \lr{X \textasciitilde{} N(0, \ensuremath{\sigma}\textsuperscript{2})}. توان: \lr{\mbox{Y = X\textsuperscript{2}}}.

برای \lr{\mbox{y \textgreater{} 0}}، جواب‌ها \lr{\mbox{x = +\ensuremath{\surd}y}} و \lr{\mbox{x = -\ensuremath{\surd}y}} هستند. \lr{\mbox{g’(x) = 2x}}.

\PSdisplay{f_Y(y) = \frac{f_X(\sqrt{y})}{2\sqrt{y}} + \frac{f_X(-\sqrt{y})}{2\sqrt{y}} = \frac{1}{\sqrt{2\pi}\sigma} \cdot \frac{e^{-y/(2\sigma^2)}}{\sqrt{y}}}

این یک \textbf{توزیع کای‌دو با 1 درجه آزادی} است (مقیاس‌شده با \lr{\mbox{\ensuremath{\sigma}\textsuperscript{2}}}).

\end{PSbox}

\begin{PSbox}[unbreakable]{ex}{مثال مهندسی: یکسوساز تمام‌موج \lr{\mbox{(Y = \textbar{}X\textbar{})}}}

ورودی \lr{X \textasciitilde{} N(0, \ensuremath{\sigma}\textsuperscript{2})}. خروجی \lr{\mbox{Y = \textbar{}X\textbar{}}}.

برای \lr{\mbox{y \textgreater{} 0}}: \lr{\mbox{x = +y}} و \lr{\mbox{x = -y}} جواب هستند. \lr{\mbox{\textbar{}g’(x)\textbar{} = 1}}.

\PSdisplay{f_Y(y) = f_X(y) + f_X(-y) = \frac{2}{\sigma\sqrt{2\pi}} e^{-y^2/(2\sigma^2)}, \quad y \geq 0}

این توزیع \textbf{نرمال تاشده} (نیم‌نرمال) است.

\end{PSbox}

\PSsection{5.5}{روش ژاکوبی}

\PSsubsection{رویه برای \lr{\mbox{Y = g(X)}}}

\begin{enumerate}
\def\labelenumi{\arabic{enumi}.}
\tightlist
\item
  تبدیل \lr{\mbox{y = g(x)}} را مشخص کنید
\item
  وارون را بیابید: \lr{\mbox{x = g\textsuperscript{-1}(y)}}
\item
  ژاکوبی را محاسبه کنید: \lr{J = \textbar{}dx/dy\textbar{} = \textbar{}d[g\textsuperscript{-1}(y)]/dy\textbar{}}
\item
  اعمال کنید: \lr{f\textsubscript{Y}(y) = f\textsubscript{X}(g\textsuperscript{-1}(y)) \ensuremath{\cdot} \textbar{}J\textbar{}}
\end{enumerate}

\begin{PSbox}[unbreakable]{ex}{مثال گام‌به‌گام: نمایی گاوسی}

\lr{X \textasciitilde{} N(\ensuremath{\mu}, \ensuremath{\sigma}\textsuperscript{2})}. \lr{\mbox{Y = e\textsuperscript{x}}} (تبدیل لگ‌نرمال).

\begin{enumerate}
\def\labelenumi{\arabic{enumi}.}
\tightlist
\item
  \lr{\mbox{g(x) = e\textsuperscript{x}}} (یکنوای صعودی)
\item
  \lr{\mbox{x = ln(y)}}، معتبر برای \lr{\mbox{y \textgreater{} 0}}
\item
  \lr{\mbox{J = \textbar{}dx/dy\textbar{} = 1/y}}
\item
  \lr{f\textsubscript{Y}(y) = f\textsubscript{X}(ln y) \ensuremath{\cdot} (1/y) = (1/(y\ensuremath{\sigma}\ensuremath{\surd}(2\ensuremath{\pi}))) \ensuremath{\cdot} exp(-(ln y - \ensuremath{\mu})\textsuperscript{2}/(2\ensuremath{\sigma}\textsuperscript{2}))}
\end{enumerate}

این همان \textbf{توزیع لگ‌نرمال} است \lr{—} بسیاری از کمیت‌های مهندسی را مدل می‌کند (سایه‌اندازی در بی‌سیم، قیمت سهام، عمر قطعات همراه با فرسودگی).

\end{PSbox}

\PSsection{5.6}{کاربرد مهندسی: رایلی از گاوسی}

\PSsubsection{استخراج توزیع رایلی}

در مخابرات، سیگنال دریافتی مؤلفه‌های هم‌فاز \lr{\mbox{(I)}} و مربعی \lr{\mbox{(Q)}} دارد:

\begin{itemize}
\tightlist
\item
  \lr{I \textasciitilde{} N(0, \ensuremath{\sigma}\textsuperscript{2})}، \lr{Q \textasciitilde{} N(0, \ensuremath{\sigma}\textsuperscript{2})}، مستقل
\end{itemize}

پوش: \lr{R = \ensuremath{\surd}(I\textsuperscript{2} + Q\textsuperscript{2})}

با استفاده از تبدیل \lr{\mbox{(I, Q) \ensuremath{\to} (R, \ensuremath{\theta})}} در مختصات قطبی:

\begin{itemize}
\tightlist
\item
  \lr{\mbox{I = R\ensuremath{\cdot}cos(\ensuremath{\theta})}}، \lr{\mbox{Q = R\ensuremath{\cdot}sin(\ensuremath{\theta})}}
\item
  ژاکوبی: \lr{\textbar{}\ensuremath{\partial}(I,Q)/\ensuremath{\partial}(R,\ensuremath{\theta})\textbar{} = R}
\end{itemize}

\lr{\mbox{PDF}} مشترک \lr{\mbox{(R, \ensuremath{\theta})}}:\PSbr{}\lr{f\textsubscript{R,\ensuremath{\Theta}}(r,\ensuremath{\theta}) = f\textsubscript{I,Q}(r\ensuremath{\cdot}cos\ensuremath{\theta}, r\ensuremath{\cdot}sin\ensuremath{\theta}) \ensuremath{\cdot} r = (r/(2\ensuremath{\pi}\ensuremath{\sigma}\textsuperscript{2})) \ensuremath{\cdot} e\textsuperscript{-r\textsuperscript{2}/(2\ensuremath{\sigma}\textsuperscript{2})}}

حاشیه‌سازی روی \lr{\mbox{\ensuremath{\theta} \ensuremath{\in} [0, 2\ensuremath{\pi})}}:

\PSdisplay{f_R(r) = \frac{r}{\sigma^2} e^{-r^2/(2\sigma^2)}, \quad r \geq 0}

این همان \textbf{توزیع رایلی} است \lr{—} بنیادی برای مدل‌سازی کانال محوشدگی بی‌سیم.

\PSsection{5.7}{مثال‌های \lr{\mbox{MATLAB}}}

\begin{PSbox}[unbreakable]{ex}{مثال 1: بررسی تبدیل خطی}

\tcbset{PScodemode/.style={unbreakable}}

\begin{Shaded}
\begin{Highlighting}[]
\CommentTok{\%\% Linear Transformation: Amplifier Output}
\VariableTok{sigma\_in} \OperatorTok{=} \FloatTok{0.1}\OperatorTok{;}  \VariableTok{G} \OperatorTok{=} \FloatTok{5}\OperatorTok{;}  \VariableTok{offset} \OperatorTok{=} \FloatTok{2}\OperatorTok{;}
\VariableTok{N} \OperatorTok{=} \FloatTok{100000}\OperatorTok{;}

\VariableTok{X} \OperatorTok{=} \VariableTok{sigma\_in} \OperatorTok{*} \VariableTok{randn}\NormalTok{(}\FloatTok{1}\OperatorTok{,} \VariableTok{N}\NormalTok{)}\OperatorTok{;}   \CommentTok{\% Input noise}
\VariableTok{Y} \OperatorTok{=} \VariableTok{G}\OperatorTok{*}\VariableTok{X} \OperatorTok{+} \VariableTok{offset}\OperatorTok{;}              \CommentTok{\% Output}

\VariableTok{figure}\OperatorTok{;}
\VariableTok{subplot}\NormalTok{(}\FloatTok{1}\OperatorTok{,}\FloatTok{2}\OperatorTok{,}\FloatTok{1}\NormalTok{)}\OperatorTok{;}
\VariableTok{histogram}\NormalTok{(}\VariableTok{X}\OperatorTok{,} \FloatTok{80}\OperatorTok{,} \SpecialStringTok{\textquotesingle{}Normalization\textquotesingle{}}\OperatorTok{,} \SpecialStringTok{\textquotesingle{}pdf\textquotesingle{}}\NormalTok{)}\OperatorTok{;} \VariableTok{hold} \VariableTok{on}\OperatorTok{;}
\VariableTok{x} \OperatorTok{=} \VariableTok{linspace}\NormalTok{(}\OperatorTok{{-}}\FloatTok{0.5}\OperatorTok{,} \FloatTok{0.5}\OperatorTok{,} \FloatTok{200}\NormalTok{)}\OperatorTok{;}
\VariableTok{plot}\NormalTok{(}\VariableTok{x}\OperatorTok{,} \VariableTok{normpdf}\NormalTok{(}\VariableTok{x}\OperatorTok{,} \FloatTok{0}\OperatorTok{,} \VariableTok{sigma\_in}\NormalTok{)}\OperatorTok{,} \SpecialStringTok{\textquotesingle{}r\textquotesingle{}}\OperatorTok{,} \SpecialStringTok{\textquotesingle{}LineWidth\textquotesingle{}}\OperatorTok{,} \FloatTok{2}\NormalTok{)}\OperatorTok{;}
\VariableTok{title}\NormalTok{(}\SpecialStringTok{\textquotesingle{}Input X \textasciitilde{} N(0, 0.01)\textquotesingle{}}\NormalTok{)}\OperatorTok{;} \VariableTok{xlabel}\NormalTok{(}\SpecialStringTok{\textquotesingle{}Voltage (V)\textquotesingle{}}\NormalTok{)}\OperatorTok{;}

\VariableTok{subplot}\NormalTok{(}\FloatTok{1}\OperatorTok{,}\FloatTok{2}\OperatorTok{,}\FloatTok{2}\NormalTok{)}\OperatorTok{;}
\VariableTok{histogram}\NormalTok{(}\VariableTok{Y}\OperatorTok{,} \FloatTok{80}\OperatorTok{,} \SpecialStringTok{\textquotesingle{}Normalization\textquotesingle{}}\OperatorTok{,} \SpecialStringTok{\textquotesingle{}pdf\textquotesingle{}}\NormalTok{)}\OperatorTok{;} \VariableTok{hold} \VariableTok{on}\OperatorTok{;}
\VariableTok{y} \OperatorTok{=} \VariableTok{linspace}\NormalTok{(}\FloatTok{0}\OperatorTok{,} \FloatTok{4}\OperatorTok{,} \FloatTok{200}\NormalTok{)}\OperatorTok{;}
\VariableTok{plot}\NormalTok{(}\VariableTok{y}\OperatorTok{,} \VariableTok{normpdf}\NormalTok{(}\VariableTok{y}\OperatorTok{,} \VariableTok{offset}\OperatorTok{,} \VariableTok{G}\OperatorTok{*}\VariableTok{sigma\_in}\NormalTok{)}\OperatorTok{,} \SpecialStringTok{\textquotesingle{}r\textquotesingle{}}\OperatorTok{,} \SpecialStringTok{\textquotesingle{}LineWidth\textquotesingle{}}\OperatorTok{,} \FloatTok{2}\NormalTok{)}\OperatorTok{;}
\VariableTok{title}\NormalTok{(}\SpecialStringTok{\textquotesingle{}Output Y = 5X + 2\textquotesingle{}}\NormalTok{)}\OperatorTok{;} \VariableTok{xlabel}\NormalTok{(}\SpecialStringTok{\textquotesingle{}Voltage (V)\textquotesingle{}}\NormalTok{)}\OperatorTok{;}
\end{Highlighting}
\end{Shaded}

\end{PSbox}

\begin{PSbox}[unbreakable]{ex}{مثال 2: توان از ولتاژ گاوسی}

\tcbset{PScodemode/.style={unbreakable}}

\begin{Shaded}
\begin{Highlighting}[]
\CommentTok{\%\% Nonlinear: Y = X\^{}2 (instantaneous power)}
\VariableTok{sigma} \OperatorTok{=} \FloatTok{1}\OperatorTok{;} \VariableTok{N} \OperatorTok{=} \FloatTok{200000}\OperatorTok{;}
\VariableTok{X} \OperatorTok{=} \VariableTok{sigma} \OperatorTok{*} \VariableTok{randn}\NormalTok{(}\FloatTok{1}\OperatorTok{,} \VariableTok{N}\NormalTok{)}\OperatorTok{;}
\VariableTok{Y} \OperatorTok{=} \VariableTok{X}\OperatorTok{.\^{}}\FloatTok{2}\OperatorTok{;}

\VariableTok{figure}\OperatorTok{;}
\VariableTok{histogram}\NormalTok{(}\VariableTok{Y}\OperatorTok{,} \FloatTok{150}\OperatorTok{,} \SpecialStringTok{\textquotesingle{}Normalization\textquotesingle{}}\OperatorTok{,} \SpecialStringTok{\textquotesingle{}pdf\textquotesingle{}}\OperatorTok{,} \SpecialStringTok{\textquotesingle{}BinLimits\textquotesingle{}}\OperatorTok{,}\NormalTok{ [}\FloatTok{0}\OperatorTok{,} \FloatTok{8}\NormalTok{])}\OperatorTok{;}
\VariableTok{hold} \VariableTok{on}\OperatorTok{;}
\VariableTok{y} \OperatorTok{=} \VariableTok{linspace}\NormalTok{(}\FloatTok{0.01}\OperatorTok{,} \FloatTok{8}\OperatorTok{,} \FloatTok{300}\NormalTok{)}\OperatorTok{;}
\VariableTok{f\_theory} \OperatorTok{=} \VariableTok{chi2pdf}\NormalTok{(}\VariableTok{y}\OperatorTok{/}\VariableTok{sigma}\OperatorTok{\^{}}\FloatTok{2}\OperatorTok{,} \FloatTok{1}\NormalTok{) }\OperatorTok{/} \VariableTok{sigma}\OperatorTok{\^{}}\FloatTok{2}\OperatorTok{;}
\VariableTok{plot}\NormalTok{(}\VariableTok{y}\OperatorTok{,} \VariableTok{f\_theory}\OperatorTok{,} \SpecialStringTok{\textquotesingle{}r\textquotesingle{}}\OperatorTok{,} \SpecialStringTok{\textquotesingle{}LineWidth\textquotesingle{}}\OperatorTok{,} \FloatTok{2}\NormalTok{)}\OperatorTok{;}
\VariableTok{xlabel}\NormalTok{(}\SpecialStringTok{\textquotesingle{}Power (V²)\textquotesingle{}}\NormalTok{)}\OperatorTok{;} \VariableTok{ylabel}\NormalTok{(}\SpecialStringTok{\textquotesingle{}PDF\textquotesingle{}}\NormalTok{)}\OperatorTok{;}
\VariableTok{title}\NormalTok{(}\SpecialStringTok{\textquotesingle{}Distribution of Power Y = X²\textquotesingle{}}\NormalTok{)}\OperatorTok{;}
\VariableTok{legend}\NormalTok{(}\SpecialStringTok{\textquotesingle{}Simulation\textquotesingle{}}\OperatorTok{,} \SpecialStringTok{\textquotesingle{}Theory (scaled χ²₁)\textquotesingle{}}\NormalTok{)}\OperatorTok{;}
\end{Highlighting}
\end{Shaded}

\end{PSbox}

\begin{PSbox}[unbreakable]{ex}{مثال 3: \lr{\mbox{SNR}} بر حسب \lr{\mbox{dB}} از \lr{\mbox{SNR}} خطی}

\tcbset{PScodemode/.style={unbreakable}}

\begin{Shaded}
\begin{Highlighting}[]
\CommentTok{\%\% SNR transformation: Y = 10*log10(X)}
\CommentTok{\% X \textasciitilde{} Exponential(1) (Rayleigh fading power, normalized)}
\VariableTok{N} \OperatorTok{=} \FloatTok{200000}\OperatorTok{;}
\VariableTok{X} \OperatorTok{=} \VariableTok{exprnd}\NormalTok{(}\FloatTok{1}\OperatorTok{,} \FloatTok{1}\OperatorTok{,} \VariableTok{N}\NormalTok{)}\OperatorTok{;}
\VariableTok{Y} \OperatorTok{=} \FloatTok{10}\OperatorTok{*}\VariableTok{log10}\NormalTok{(}\VariableTok{X}\NormalTok{)}\OperatorTok{;}   \CommentTok{\% SNR in dB}

\VariableTok{figure}\OperatorTok{;}
\VariableTok{histogram}\NormalTok{(}\VariableTok{Y}\OperatorTok{,} \FloatTok{150}\OperatorTok{,} \SpecialStringTok{\textquotesingle{}Normalization\textquotesingle{}}\OperatorTok{,} \SpecialStringTok{\textquotesingle{}pdf\textquotesingle{}}\NormalTok{)}\OperatorTok{;}
\VariableTok{hold} \VariableTok{on}\OperatorTok{;}
\VariableTok{y} \OperatorTok{=} \VariableTok{linspace}\NormalTok{(}\OperatorTok{{-}}\FloatTok{30}\OperatorTok{,} \FloatTok{15}\OperatorTok{,} \FloatTok{300}\NormalTok{)}\OperatorTok{;}
\CommentTok{\% Theoretical PDF of Y = 10*log10(X) where X \textasciitilde{} Exp(1)}
\VariableTok{f\_Y} \OperatorTok{=}\NormalTok{ (}\VariableTok{log}\NormalTok{(}\FloatTok{10}\NormalTok{)}\OperatorTok{/}\FloatTok{10}\NormalTok{) }\OperatorTok{.*} \FloatTok{10}\OperatorTok{.\^{}}\NormalTok{(}\VariableTok{y}\OperatorTok{/}\FloatTok{10}\NormalTok{) }\OperatorTok{.*} \VariableTok{exp}\NormalTok{(}\OperatorTok{{-}}\FloatTok{10}\OperatorTok{.\^{}}\NormalTok{(}\VariableTok{y}\OperatorTok{/}\FloatTok{10}\NormalTok{))}\OperatorTok{;}
\VariableTok{plot}\NormalTok{(}\VariableTok{y}\OperatorTok{,} \VariableTok{f\_Y}\OperatorTok{,} \SpecialStringTok{\textquotesingle{}r\textquotesingle{}}\OperatorTok{,} \SpecialStringTok{\textquotesingle{}LineWidth\textquotesingle{}}\OperatorTok{,} \FloatTok{2}\NormalTok{)}\OperatorTok{;}
\VariableTok{xlabel}\NormalTok{(}\SpecialStringTok{\textquotesingle{}SNR (dB)\textquotesingle{}}\NormalTok{)}\OperatorTok{;} \VariableTok{ylabel}\NormalTok{(}\SpecialStringTok{\textquotesingle{}PDF\textquotesingle{}}\NormalTok{)}\OperatorTok{;}
\VariableTok{title}\NormalTok{(}\SpecialStringTok{\textquotesingle{}Distribution of SNR in dB (Rayleigh Fading)\textquotesingle{}}\NormalTok{)}\OperatorTok{;}
\VariableTok{legend}\NormalTok{(}\SpecialStringTok{\textquotesingle{}Simulation\textquotesingle{}}\OperatorTok{,} \SpecialStringTok{\textquotesingle{}Theory\textquotesingle{}}\NormalTok{)}\OperatorTok{;}
\end{Highlighting}
\end{Shaded}

\end{PSbox}

\begin{PSbox}[unbreakable]{ex}{مثال 4: پوش رایلی از مؤلفه‌های گاوسی}

\tcbset{PScodemode/.style={unbreakable}}

\begin{Shaded}
\begin{Highlighting}[]
\CommentTok{\%\% Envelope of Complex Gaussian}
\VariableTok{sigma} \OperatorTok{=} \FloatTok{1}\OperatorTok{;} \VariableTok{N} \OperatorTok{=} \FloatTok{100000}\OperatorTok{;}
\VariableTok{I} \OperatorTok{=} \VariableTok{sigma} \OperatorTok{*} \VariableTok{randn}\NormalTok{(}\FloatTok{1}\OperatorTok{,} \VariableTok{N}\NormalTok{)}\OperatorTok{;}
\VariableTok{Q} \OperatorTok{=} \VariableTok{sigma} \OperatorTok{*} \VariableTok{randn}\NormalTok{(}\FloatTok{1}\OperatorTok{,} \VariableTok{N}\NormalTok{)}\OperatorTok{;}
\VariableTok{R} \OperatorTok{=} \VariableTok{sqrt}\NormalTok{(}\VariableTok{I}\OperatorTok{.\^{}}\FloatTok{2} \OperatorTok{+} \VariableTok{Q}\OperatorTok{.\^{}}\FloatTok{2}\NormalTok{)}\OperatorTok{;}

\VariableTok{figure}\OperatorTok{;}
\VariableTok{histogram}\NormalTok{(}\VariableTok{R}\OperatorTok{,} \FloatTok{100}\OperatorTok{,} \SpecialStringTok{\textquotesingle{}Normalization\textquotesingle{}}\OperatorTok{,} \SpecialStringTok{\textquotesingle{}pdf\textquotesingle{}}\NormalTok{)}\OperatorTok{;}
\VariableTok{hold} \VariableTok{on}\OperatorTok{;}
\VariableTok{r} \OperatorTok{=} \VariableTok{linspace}\NormalTok{(}\FloatTok{0}\OperatorTok{,} \FloatTok{5}\OperatorTok{,} \FloatTok{200}\NormalTok{)}\OperatorTok{;}
\VariableTok{plot}\NormalTok{(}\VariableTok{r}\OperatorTok{,} \VariableTok{raylpdf}\NormalTok{(}\VariableTok{r}\OperatorTok{,} \VariableTok{sigma}\NormalTok{)}\OperatorTok{,} \SpecialStringTok{\textquotesingle{}r\textquotesingle{}}\OperatorTok{,} \SpecialStringTok{\textquotesingle{}LineWidth\textquotesingle{}}\OperatorTok{,} \FloatTok{2}\NormalTok{)}\OperatorTok{;}
\VariableTok{xlabel}\NormalTok{(}\SpecialStringTok{\textquotesingle{}Envelope |h|\textquotesingle{}}\NormalTok{)}\OperatorTok{;} \VariableTok{ylabel}\NormalTok{(}\SpecialStringTok{\textquotesingle{}PDF\textquotesingle{}}\NormalTok{)}\OperatorTok{;}
\VariableTok{title}\NormalTok{(}\SpecialStringTok{\textquotesingle{}Rayleigh Distribution from I/Q Components\textquotesingle{}}\NormalTok{)}\OperatorTok{;}
\VariableTok{legend}\NormalTok{(}\SpecialStringTok{\textquotesingle{}Simulation\textquotesingle{}}\OperatorTok{,} \SpecialStringTok{\textquotesingle{}Theory\textquotesingle{}}\NormalTok{)}\OperatorTok{;}
\VariableTok{fprintf}\NormalTok{(}\SpecialStringTok{\textquotesingle{}Mean: Sim=\%.3f, Theory=\%.3f\textbackslash{}n\textquotesingle{}}\OperatorTok{,} \VariableTok{mean}\NormalTok{(}\VariableTok{R}\NormalTok{)}\OperatorTok{,} \VariableTok{sigma}\OperatorTok{*}\VariableTok{sqrt}\NormalTok{(}\VariableTok{pi}\OperatorTok{/}\FloatTok{2}\NormalTok{))}\OperatorTok{;}
\end{Highlighting}
\end{Shaded}

\end{PSbox}

\begin{PSbox}[unbreakable]{ex}{مثال 5: لگ‌نرمال از گاوسی}

\tcbset{PScodemode/.style={unbreakable}}

\begin{Shaded}
\begin{Highlighting}[]
\CommentTok{\%\% Log{-}Normal: Y = exp(X) where X \textasciitilde{} N(mu, sigma\^{}2)}
\VariableTok{mu} \OperatorTok{=} \FloatTok{0}\OperatorTok{;} \VariableTok{sigma\_x} \OperatorTok{=} \FloatTok{0.5}\OperatorTok{;} \VariableTok{N} \OperatorTok{=} \FloatTok{100000}\OperatorTok{;}
\VariableTok{X} \OperatorTok{=} \VariableTok{mu} \OperatorTok{+} \VariableTok{sigma\_x}\OperatorTok{*}\VariableTok{randn}\NormalTok{(}\FloatTok{1}\OperatorTok{,} \VariableTok{N}\NormalTok{)}\OperatorTok{;}
\VariableTok{Y} \OperatorTok{=} \VariableTok{exp}\NormalTok{(}\VariableTok{X}\NormalTok{)}\OperatorTok{;}

\VariableTok{figure}\OperatorTok{;}
\VariableTok{histogram}\NormalTok{(}\VariableTok{Y}\OperatorTok{,} \FloatTok{150}\OperatorTok{,} \SpecialStringTok{\textquotesingle{}Normalization\textquotesingle{}}\OperatorTok{,} \SpecialStringTok{\textquotesingle{}pdf\textquotesingle{}}\OperatorTok{,} \SpecialStringTok{\textquotesingle{}BinLimits\textquotesingle{}}\OperatorTok{,}\NormalTok{ [}\FloatTok{0}\OperatorTok{,} \FloatTok{5}\NormalTok{])}\OperatorTok{;}
\VariableTok{hold} \VariableTok{on}\OperatorTok{;}
\VariableTok{y} \OperatorTok{=} \VariableTok{linspace}\NormalTok{(}\FloatTok{0.01}\OperatorTok{,} \FloatTok{5}\OperatorTok{,} \FloatTok{300}\NormalTok{)}\OperatorTok{;}
\VariableTok{f\_lognorm} \OperatorTok{=} \VariableTok{lognpdf}\NormalTok{(}\VariableTok{y}\OperatorTok{,} \VariableTok{mu}\OperatorTok{,} \VariableTok{sigma\_x}\NormalTok{)}\OperatorTok{;}
\VariableTok{plot}\NormalTok{(}\VariableTok{y}\OperatorTok{,} \VariableTok{f\_lognorm}\OperatorTok{,} \SpecialStringTok{\textquotesingle{}r\textquotesingle{}}\OperatorTok{,} \SpecialStringTok{\textquotesingle{}LineWidth\textquotesingle{}}\OperatorTok{,} \FloatTok{2}\NormalTok{)}\OperatorTok{;}
\VariableTok{xlabel}\NormalTok{(}\SpecialStringTok{\textquotesingle{}Y = e\^{}X\textquotesingle{}}\NormalTok{)}\OperatorTok{;} \VariableTok{ylabel}\NormalTok{(}\SpecialStringTok{\textquotesingle{}PDF\textquotesingle{}}\NormalTok{)}\OperatorTok{;}
\VariableTok{title}\NormalTok{(}\SpecialStringTok{\textquotesingle{}Log{-}Normal Distribution\textquotesingle{}}\NormalTok{)}\OperatorTok{;}
\VariableTok{legend}\NormalTok{(}\SpecialStringTok{\textquotesingle{}Simulation\textquotesingle{}}\OperatorTok{,} \SpecialStringTok{\textquotesingle{}Theory\textquotesingle{}}\NormalTok{)}\OperatorTok{;}
\end{Highlighting}
\end{Shaded}

\end{PSbox}

\PSsection{5.8}{مسایل تمرینی}

\begin{PSbox}[unbreakable]{prob}{مسیله 1}

سیگنال \lr{\mbox{X \textasciitilde{} Uniform[0, 1]}}. خروجی \lr{\mbox{Y = -2\ensuremath{\cdot}ln(X)}}. \lr{\mbox{PDF}} مربوط به \lr{\mbox{Y}} را بیابید و توزیع را تشخیص دهید.

\PSsolution{} روش \lr{\mbox{CDF}}: \lr{F\textsubscript{Y}(y) = P(-2ln(X) \ensuremath{\le} y) = P(X \ensuremath{\ge} e\textsuperscript{-y/2}) = 1 - e\textsuperscript{-y/2}} برای \lr{\mbox{y \ensuremath{\ge} 0}}.\PSbr{}\lr{f\textsubscript{Y}(y) = (1/2)e\textsuperscript{-y/2} \ensuremath{\to} Y \textasciitilde{} Exponential(\ensuremath{\beta} = 2)}. (این همان روش \lr{\mbox{CDF}} وارون برای تولید توزیع نمایی است!)

\end{PSbox}

\begin{PSbox}[unbreakable]{prob}{مسیله 2}

ولتاژ نویز \lr{X \textasciitilde{} N(0, \ensuremath{\sigma}\textsuperscript{2})} با \lr{\mbox{\ensuremath{\sigma} = 2V}}. توان تلف‌شده در \lr{\mbox{R = 50\ensuremath{\Omega}}}: \lr{\mbox{P = X\textsuperscript{2}/R}}.

\begin{enumerate}
\def\labelenumi{(\PSalph{enumi})}
\tightlist
\item
  \lr{\mbox{E[P]}} را بیابید. (ب) \lr{\mbox{PDF}} مربوط به \lr{\mbox{P}} را بیابید. (ج) \lr{\mbox{P(P \textgreater{} 0.2W)}} را بیابید.
\end{enumerate}

\PSsolution{} 

\begin{enumerate}
\def\labelenumi{(\PSalph{enumi})}
\tightlist
\item
  \lr{E[P] = E[X\textsuperscript{2}]/R = \ensuremath{\sigma}\textsuperscript{2}/R = 4/50 = 0.08W}
\item
  \lr{\mbox{P = X\textsuperscript{2}/50}}. فرض کنید \lr{W = X\textsuperscript{2} \textasciitilde{} \ensuremath{\sigma}\textsuperscript{2}\ensuremath{\cdot}\ensuremath{\chi}\textsuperscript{2}(1)}. آنگاه \lr{\mbox{P = W/50}}. \lr{f\textsubscript{P}(p) = 50\ensuremath{\cdot}f\textsubscript{W}(50p) = (50/(2\ensuremath{\sigma}\textsuperscript{2}))\ensuremath{\cdot}(50p/\ensuremath{\sigma}\textsuperscript{2})\textsuperscript{-1/2}\ensuremath{\cdot}e\textsuperscript{-50p/(2\ensuremath{\sigma}\textsuperscript{2})}} برای \lr{\mbox{p \textgreater{} 0}}.
\item
  از \lr{\mbox{MATLAB}} استفاده کنید: \lr{\PScode{1 - chi2cdf(0.2*50/4, 1)} = 1 - chi2cdf(2.5, 1) \ensuremath{\approx} 0.114}
\end{enumerate}

\end{PSbox}

\begin{PSbox}[unbreakable]{prob}{مسیله 3}

\lr{\mbox{X \textasciitilde{} Exponential(\ensuremath{\lambda} = 1)}}. \lr{\mbox{Y = \ensuremath{\surd}X}}. مقدار \lr{\mbox{f\textsubscript{Y}(y)}} را بیابید.

\PSsolution{} \lr{g(x) = \ensuremath{\surd}x \ensuremath{\to} x = y\textsuperscript{2}}، \lr{\mbox{dx/dy = 2y}}.\PSbr{}\lr{f\textsubscript{Y}(y) = f\textsubscript{X}(y\textsuperscript{2})\ensuremath{\cdot}\textbar{}2y\textbar{} = e\textsuperscript{-y\textsuperscript{2}}\ensuremath{\cdot}2y} برای \lr{\mbox{y \ensuremath{\ge} 0}}. این یک توزیع رایلی با \lr{\mbox{\ensuremath{\sigma}\textsuperscript{2} = 1/2}} است.

\end{PSbox}

\begin{PSbox}[unbreakable]{prob}{مسیله 4}

فاصله \lr{\mbox{D \textasciitilde{} Uniform[1, 10] km}}. تلفات مسیر: \lr{\mbox{L = 20\ensuremath{\cdot}log\textsubscript{10}(D) dB}}. مقدار \lr{\mbox{E[L]}} و \lr{\mbox{PDF}} مربوط به \lr{\mbox{L}} را بیابید.

\PSsolution{} \lr{E[L] = E[20\ensuremath{\cdot}log\textsubscript{10}(D)] = \ensuremath{\int}\textsubscript{1}\textsuperscript{10} 20\ensuremath{\cdot}log\textsubscript{10}(d)\ensuremath{\cdot}(1/9) dd = (20/9)\ensuremath{\cdot}\ensuremath{\int}\textsubscript{1}\textsuperscript{10} log\textsubscript{10}(d) dd}\PSbr{}\lr{= (20/9)\ensuremath{\cdot}[d\ensuremath{\cdot}log\textsubscript{10}(d) - d/ln(10)]\textsubscript{1}\textsuperscript{10} = (20/9)\ensuremath{\cdot}[10 - 10/ln10 + 1/ln10] = (20/9)\ensuremath{\cdot}(10 - 9/ln10) \ensuremath{\approx} 13.55 dB}

\lr{\mbox{PDF}}: \lr{\mbox{L}} از 0 تا \lr{\mbox{20 dB}} تغییر می‌کند. \lr{\mbox{g\textsuperscript{-1}(l) = 10\textsuperscript{l/20}}}، \lr{\textbar{}dg\textsuperscript{-1}/dl\textbar{} = (ln10/20)\ensuremath{\cdot}10\textsuperscript{l/20}}.\PSbr{}\lr{f\textsubscript{L}(l) = (1/9)\ensuremath{\cdot}(ln10/20)\ensuremath{\cdot}10\textsuperscript{l/20}} برای \lr{\mbox{0 \ensuremath{\le} l \ensuremath{\le} 20}}.

\end{PSbox}

\begin{PSbox}[unbreakable]{prob}{مسیله 5}

کد \lr{\mbox{MATLAB}} را برای بررسی مسیله 2 با شبیه‌سازی بنویسید.

\PSsolution{} 

\tcbset{PScodemode/.style={unbreakable}}

\begin{Shaded}
\begin{Highlighting}[]
\VariableTok{sigma} \OperatorTok{=} \FloatTok{2}\OperatorTok{;} \VariableTok{R} \OperatorTok{=} \FloatTok{50}\OperatorTok{;} \VariableTok{N} \OperatorTok{=} \FloatTok{500000}\OperatorTok{;}
\VariableTok{X} \OperatorTok{=} \VariableTok{sigma}\OperatorTok{*}\VariableTok{randn}\NormalTok{(}\FloatTok{1}\OperatorTok{,}\VariableTok{N}\NormalTok{)}\OperatorTok{;}
\VariableTok{P} \OperatorTok{=} \VariableTok{X}\OperatorTok{.\^{}}\FloatTok{2} \OperatorTok{/} \VariableTok{R}\OperatorTok{;}
\VariableTok{fprintf}\NormalTok{(}\SpecialStringTok{\textquotesingle{}E[P]: Theory=\%.4f W, Sim=\%.4f W\textbackslash{}n\textquotesingle{}}\OperatorTok{,} \VariableTok{sigma}\OperatorTok{\^{}}\FloatTok{2}\OperatorTok{/}\VariableTok{R}\OperatorTok{,} \VariableTok{mean}\NormalTok{(}\VariableTok{P}\NormalTok{))}\OperatorTok{;}
\VariableTok{fprintf}\NormalTok{(}\SpecialStringTok{\textquotesingle{}P(P\textgreater{}0.2): Theory=\%.4f, Sim=\%.4f\textbackslash{}n\textquotesingle{}}\OperatorTok{,} \FloatTok{1}\OperatorTok{{-}}\VariableTok{chi2cdf}\NormalTok{(}\FloatTok{0.2}\OperatorTok{*}\VariableTok{R}\OperatorTok{/}\VariableTok{sigma}\OperatorTok{\^{}}\FloatTok{2}\OperatorTok{,}\FloatTok{1}\NormalTok{)}\OperatorTok{,} \VariableTok{mean}\NormalTok{(}\VariableTok{P}\OperatorTok{\textgreater{}}\FloatTok{0.2}\NormalTok{))}\OperatorTok{;}
\end{Highlighting}
\end{Shaded}

\emph{ماژول بعدی: {ماژول 6 \lr{—} دو متغیر تصادفی}}

\end{PSbox}
